\documentclass[12pt,letterpaper]{article}
\usepackage[utf8]{inputenc}
\usepackage[spanish]{babel}
\usepackage{amsmath}
\usepackage{amsfonts}
\usepackage{amssymb}
\usepackage[margin=2.5cm]{geometry}

\begin{document}

\begin{center}
    \Large \textbf{Evaluación de Examen 2: Unidad II} \\
    \large Física para Ciencias de la Salud (005-1713) \\
    \vspace{0.5cm}
    \textbf{Estudiante:} Daniela Escribano \quad \textbf{C.I.:} 32.031.568
    \vspace{0.3cm} \\
    \large \textbf{Calificación Final: 9/10 puntos}
\end{center}

\vspace{0.5cm}

\section*{Análisis de Resultados}

\subsection*{1. Cinemática (Aceleración media)}
\textbf{Procedimiento y resultado:} Extrae los datos correctamente y aplica la fórmula de aceleración media ($a = \frac{v_f - v_i}{t}$). La sustitución y cálculo son precisos, obteniendo el valor exacto de $4 \text{ m/s}^2$ con sus unidades correctas en el Sistema Internacional. \\
\textbf{Veredicto:} \textbf{Correcto} (2/2 pts).

\subsection*{2. Dinámica (Leyes de Newton)}
\textbf{Procedimiento y resultado:} Plantea la Segunda Ley de Newton ($F = m \cdot a$) y despeja de forma analítica e impecable la aceleración ($a = \frac{F}{m}$). Sustituye los valores para obtener correctamente $6 \text{ m/s}^2$. \\
\textbf{Veredicto:} \textbf{Correcto} (2/2 pts).

\subsection*{3. Trabajo Mecánico}
\textbf{Procedimiento y resultado:} Extrae los datos y destaca acertadamente de forma teórica que la fuerza y el desplazamiento tienen la misma dirección, validando el uso directo de la fórmula. Aplica la multiplicación para el trabajo ($W = F \cdot d$) y obtiene $320$, indicando explícitamente el resultado final en Joules. \\
\textbf{Veredicto:} \textbf{Correcto} (2/2 pts).

\subsection*{4. Energía Potencial}
\textbf{Procedimiento y resultado:} Plantea la fórmula de la energía potencial gravitatoria ($E_p = m \cdot g \cdot h$). Sustituye los valores correspondientes incluyendo la constante de gravedad de $9,8 \text{ m/s}^2$ y calcula exactamente $17,64 \text{ J}$. \\
\textbf{Veredicto:} \textbf{Correcto} (2/2 pts).

\subsection*{5. Potencia Mecánica}
\textbf{Procedimiento y resultado:} La estudiante extrae los datos de forma correcta e identifica la fórmula a utilizar ($P = \frac{W}{t}$). Sin embargo, el desarrollo del problema queda inconcluso en este punto, omitiendo la sustitución numérica de los $75 \text{ J}$ y $60 \text{ s}$, lo que le impide llegar al resultado final ($1,25 \text{ W}$). \\
\textbf{Veredicto:} \textbf{Incompleto} (Se otorgan puntos parciales por la identificación correcta de los datos y el modelo físico) (1/2 pts).

\vspace{0.5cm}
\section*{Comentarios Generales y Sugerencias}
¡Un examen excelente! La estudiante demuestra un dominio sobresaliente de los conceptos de la Unidad II y goza de una caligrafía y orden de resolución dignos de destacar.
\begin{itemize}
    \item La metodología aplicada (datos $\rightarrow$ despejes $\rightarrow$ sustitución $\rightarrow$ enmarcado del resultado) es la idónea y se recomienda fuertemente mantenerla en evaluaciones futuras.
    \item Como un pequeño detalle formal, se recuerda que el símbolo oficial en el Sistema Internacional para kilogramos es minúscula ($\text{kg}$), en lugar de $\text{Kg}$.
    \item Se aconseja vigilar el manejo del tiempo durante las evaluaciones para poder concluir todos los ejercicios satisfactoriamente, evitando que problemas bien encaminados (como la pregunta 5) queden sin resolver por completo.
\end{itemize}

\end{document}